%% Reader guidance for The Methods of Reason.
%% The obsolete inherited title, preface and graduate-course scope are archival.
\frontmatter
\tableofcontents
\mainmatter

\chapter*{How to Use This Handbook}
\addcontentsline{toc}{chapter}{How to Use This Handbook}
\markboth{HOW TO USE THIS HANDBOOK}{HOW TO USE THIS HANDBOOK}

\emph{The Methods of Reason} is a reference to reasoning and logic.
It brings together deductive, inductive, probabilistic, statistical,
Bayesian, abductive, causal, decision-theoretic, strategic, and informal
methods, together with empirical research on how people reason.
The organization supports consultation by subject and by the question
a method is intended to answer.

\section*{Finding and using a treatment}

Use the contents, chapter headings, and subject index to locate a topic.
For a reasoning task, identify what information is available, what
conclusion or decision is sought, and which assumptions can be justified.
The relevant method supplies its inputs, procedure, worked cases, and
conditions of application.

Numbered entries include definitions, conventions, methods, examples,
theorems, and other results. They connect concise statements to full
treatments and exact sources. Definitions introduce the objects and
notation before they are used. Results state their hypotheses, including
any finiteness, countability, consistency, effectiveness, independence,
or choice assumptions on which they depend.

Examples show how a method operates. Counterexamples and failure cases
show what can change when a hypothesis is removed. Comparisons identify
which definitions, assumptions, and conclusions the compared approaches
share and where they differ.

\section*{Reading the evidence}

A mathematical result is supported by an argument under stated
hypotheses. An empirical finding is tied to its task, population, design,
measures, uncertainty, and limits of generalization. A normative argument
states the goals, values, or decision criteria on which its recommendation
depends. These forms of support have different roles.

The source references identify the presentation being used. Materially
different systems, editions, and formulations retain their own definitions
and qualifications. Corrections appear directly in the statements and
arguments where they are needed.

\section*{Proof-status taglines}
\index{status tagline}

Mathematical theorem-class entries use three proof-status forms:
\begin{description}
\item[proved (Annex \(P n.k\)).]
The specified annex item supplies a complete proof of the stated result.
Dependencies and any separately invoked results are identified in that
proof.
\item[sketched (Annex \(P n.k\)).]
The annex item gives the construction and identifies the substantive
steps that remain omitted.
\item[cited (source).]
The entry gives the exact statement and hypotheses and identifies the
source carrying its proof.
\end{description}
These tags describe the support supplied for a particular statement.
Empirical findings and informal arguments use the evidential explanations
appropriate to their methods.

\section*{The proof annex}
\index{proof annex}

Appendix P contains extended proofs grouped with their corresponding
subject parts. Each item identifies the handbook result it establishes,
and the statement links back to that item. Proofs already included within
an attributed treatment are read with that treatment. Sketches identify
their omissions explicitly.

\section*{Calculi, notation, and source conventions}

The calculus sheets record the language, sequent form, rules, side
conditions, negation treatment, and equality treatment of the identified
system. Their references lead to the full presentation and its results.
The symbol registry records a symbol's spelling, role, scope, precedence,
and first use.

A source presentation has its own notation scope and retains its
identified syntax. When a symbol has different roles in different
systems, each role is declared where it applies. A comparison between
systems gives the translation or other connection it uses.

\chapter*{Editorial Conventions and Evidence Policy}
\addcontentsline{toc}{chapter}{Editorial Conventions and Evidence Policy}
\markboth{EDITORIAL CONVENTIONS}{EDITORIAL CONVENTIONS}

\section*{Definitions and terminology}

Technical terms are defined before substantive use. Formal objects
receive their types and defining conditions; empirical constructs receive
the operational meaning used by the identified study.
Departures from a shared convention are stated with their scope.

For a new formal technical name, the treatment states its necessity,
definition, formal standing, distinction from existing terms, and use
in the subsequent argument. These conditions make the introduced
distinction explicit. Established source terminology remains attributed
to its particular presentation.

\section*{What different checks establish}

A calculation establishes its calculated instances. A truth table covers
the valuations it enumerates. A derivation establishes derivability in
its specified calculus, and a proof establishes its conclusion under its
hypotheses. A proof-assistant check is relative to the kernel, imported
declarations, and axioms used. Typesetting checks establish properties of
the document's production. Independent mathematical, empirical, and
editorial assessments address their respective questions.

The handbook identifies these forms of evidence at their actual scope.
Source attribution, a computed example, a proof, an empirical observation,
and a production check each support the claim for which they are relevant.
